\begin{statement}{P3K.29}{Refuted; corrected scalar-fan computation}
\par\noindent\textbf{Audit correction, 17 November 2026.} The claimed simultaneous marked delta-equivariant plumbing product does not exist for nonzero scalar winding J: P3K.61 derives the incompatible determinant equation -J=0. In the actual native C1 product, P3K.62 gives the single free pair (J,-J), contradicting the stated zero right coefficient at J=1. The integral cells, common-exterior bases, dual maps, symmetry and exact local frame-change identity remain valid in their displayed markings. The earlier statement and argument below are preserved as the located prior claim.
\medskip

For the actual index-one cusp with its three complex double curves,
use positive complex normal meridians \(m_i\) and the relative dual
basis \(m_j\cdot D_i=\delta_{ji}\). The literal common-exterior map
under the flop of \(C_0\) is
\[
 m'_0=-m_0+2m_1+2m_2,\quad m'_1=m_1,\quad m'_2=m_2,
 \qquad D'_0=-D_0,\quad D'_j=D_j+2D_0\ (j=1,2).
\]
For the scalar-parameter families and the extending cusp products
of K22--27, in a fan retaining the physical delta edge the actual
relative return has variation \(\operatorname{diag}(v,0,0)\), where
\(v=k_{\rm Jin}\) in the stated complex-meridian dual basis.
For K22--27, \(v=-3d\). After the fixed-edge flop its exact
variation matrix is
\[
 v\begin{pmatrix}1&-2&-2\\-2&4&4\\-2&4&4\end{pmatrix}.
\]
The two original nonfixed curves have integer relative derivative
coefficient \(n_{\rm rel}=0\) in K23's specified product framing.
No relative Pontryagin coefficient after the fixed-edge flop is
assigned by this statement.
\end{statement}
\begin{proof}[Preserved earlier argument; scope corrected above]
\par\medskip\noindent\textbf{1. The common exterior and its oriented bases.}\enspace

Let \(C=C_0\cup C_1\cup C_2\) be the three double curves of the
index-one cusp, in its actual smooth six-sphere. Each curve is a
complex projective line. They meet at the same two triple points,
where the actual charts have \(q=x_0x_1x_2\) and the three curves
are the coordinate axes. Put
\[
 E_C=S^6\setminus\mathring N(C).
\]
The neighbourhood is formed from small coordinate balls at those
points and complex normal tubes along the curves, with all faces
rounded equivariantly.

Here is the complete marked cellular calculation. Give each
\(C_i\) two oriented edges \(a_i,b_i\) from the first vertex
to the second and two oriented faces with boundaries
\(a_i-b_i\) and \(b_i-a_i\). In the order
\((a_0,b_0,a_1,b_1,a_2,b_2)\), its chain complex has
\[
 C_2=\mathbb Z^6\xrightarrow{D_2}C_1=\mathbb Z^6
                      \xrightarrow{D_1}C_0=\mathbb Z^2,
 \qquad
 D_1=\begin{pmatrix}-1&-1&-1&-1&-1&-1\\
                     1&1&1&1&1&1\end{pmatrix},
\]
\[
 D_2=\operatorname{diag}(B,B,B),\qquad
               B=\begin{pmatrix}1&-1\\-1&1\end{pmatrix}.
\]
The kernel of \(D_2\) is freely generated by the three
sums of paired faces, which are the positive complex
spheres. Modulo the image of \(D_2\), the two edges
\(a_i,b_i\) become the same edge with coefficient one.
The remaining differential is
\((-1,1)\) on each of these three edges. Consequently
\[
 H_2(C)=\bigoplus_i\mathbb Z[C_i],\qquad
 H_1(C)=\mathbb Z[a_1-a_0]\oplus\mathbb Z[a_2-a_0].
\]
All the identifications used unit coefficients; these
are primitive integral bases.

Triangulate the two coordinate-ball models and their
normal tubes compatibly with this cell structure.
In the oriented six-sphere, assign to each cell its
oriented transverse dual cell, with intersection
coefficient \(+1\). Passing to the complement pairs the
reduced complement chains with the cochains of \(C\),
shifted by five, and reverses the boundary incidence
in the usual dual-cell order. The map is integral:
each displayed cell is met once by its own normal
dual cell and not by the others. The cells of the
ambient six-sphere supply only the degree-zero
constant and the ambient orientation class, so the
resulting reduced complement homology is precisely
\(\widetilde H_j(E_C)=\widetilde H^{5-j}(C)\).
This gives
\[
 H_3(E_C;\mathbb Z)=\bigoplus_{i=0}^2\mathbb Z m_i,\qquad
 H_4(E_C;\mathbb Z)=\mathbb Z^2,\qquad H_2(E_C;\mathbb Z)=0.
                                                        \tag{P3K.29.1}
\]
Here \(m_i\) is the boundary of a small positively oriented complex
normal four-disc at a smooth point of \(C_i\). This fixes the
orientation by a local chart, independently of a quaternion
trivialization of a different curve's exterior. The complement is
simply connected, by relative general position. All positive
homology groups not displayed in (P3K.29.1) vanish.

Write \(D_i\in H_3(E_C,\partial E_C;\mathbb Z)\) for the
Poincare--Lefschetz dual basis:
\[
                    m_j\cdot D_i=\delta_{ji}.
                                                        \tag{P3K.29.2}
\]
For completeness, its relative duality is also integral.
Small representatives of the three normal meridians
are pairwise disjoint, so the absolute middle
intersection homomorphism is zero. The relative
duality chain map of the same triangulation identifies
\(H_3(E_C,\partial E_C)\) with \(H^3(E_C)\), whose
evaluation on the three primitive meridians is the
full \(\mathbb Z^3\). Thus (P3K.29.2) defines its entire
relative group. The boundary sequence further gives
\[
 H_j(\partial E_C)=
 \begin{cases}
  \mathbb Z^2,&j=1,4,\\
  \mathbb Z^3,&j=2,3,
 \end{cases}
\]
with the degree-three and degree-four maps to \(E_C\)
isomorphisms. In particular no division by two occurs
at a triple point.

For any chosen \(C_i\), filling the other two curve tubes identifies
its single-curve exterior \(E_i\) with \(D^3\times S^3\), in the
marking already established in K23. Its map on absolute third
homology is \(m_j\mapsto\delta_{ij}m_i\). A relative three-disc
representing the \(D^3\) factor can be chosen to miss the other two
curves: first choose its boundary normal section away from their
four specified directions at the two vertices, then apply relative
general position in the interior, since \(3+2<6\).
Its class in the common exterior is \(D_i\). Small normal
meridians of the other two curves miss this compact disc, while
the marked meridian of \(C_i\) meets it once. This is the actual
map used below to compare a common-exterior variation with K23's
single-curve variation.

\par\medskip\noindent\textbf{2. The reciprocal-section transfer chain.}\enspace

Flop \(C_0\), and denote the three new curves by \(C'_0,C'_1,C'_2\),
where the latter two are the strict transforms. The common blowup
has exceptional divisor
\[
          E=\mathbb P^1_a\times\mathbb P^1_c,\qquad
          N_E=L=\mathcal O_E(-1,-1).
                                                        \tag{P3K.29.3}
\]
This follows from the actual blowup of the zero section in
\(\mathcal O(-1)\otimes\mathbb C^2\): \(a\) is its old base
coordinate, \(c\) is the projective direction of its two normal
coordinates, and the exceptional scalar is a vector in \(L\).
The two blowdowns are the two projections of (P3K.29.3).

The strict transforms of \(C_1,C_2\) meet \(E\) in four distinct
points. The two local coordinate axes at each endpoint of \(C_0\)
give its two distinct normal directions. Each of \(C_1,C_2\)
therefore supplies two of these points. Remove four disjoint
small complex coordinate balls from \(E\), and write \(E^\circ\)
for the resulting four-manifold.

Let \(T\) denote oriented circle transfer in \(S(L)|_{E^\circ}\),
with the fibre oriented by increasing complex argument of the
exceptional scalar. At a fixed \(a\), the old complex normal
meridian is exactly \(T[c]\): it is the unit sphere of the two
normal complex coordinates before blowup. At fixed \(c\), the new
complex normal meridian is exactly \(T[a]\), by the other blowdown.
Thus both signs follow from the actual complex normal four-discs.

Choose a smooth divisor \(D\) of bidegree \((1,1)\), avoiding the
four balls, as the zero locus of a transverse section \(\sigma\)
of \(L^\vee\). On the complement of a small tube about \(D\),
the normalized reciprocal \(\sigma^{-1}/|\sigma^{-1}|\) is an
actual section of \(S(L)\). Its graph is an oriented four-chain.
Along \(D\), let \(\sigma=re^{i\phi}\) in a positive local normal
coordinate. The boundary orientation of the complement of its
tube is
\(-\operatorname{or}(D)\wedge d\phi\). The reciprocal section
has fibre argument \(-\phi\), so this boundary maps to
\(+T[D]\). At each of the four deleted balls the boundary is
minus its positive three-sphere. Consequently the graph has
the literal oriented boundary
\[
           T[a]+T[c]-\mu_{00}-\mu_{01}-\mu_{10}-\mu_{11}.
                                                        \tag{P3K.29.4}
\]
Here \(T[D]=T[a]+T[c]\) follows by the homology transfer of
\([D]=[a]+[c]\) in \(E^\circ\).

Each small section lift \(\mu\) is the positive complex normal
meridian of the strict curve meeting that point. Indeed, in the
blowup chart that strict curve is transverse to \(E\); the two
coordinates on \(E\) are its two complex normal coordinates.
Fixing a nonzero exceptional scalar identifies its small normal
three-sphere with the indicated section lift. A change of unit
frame on that three-sphere has degree zero in \(S^1\), so it
does not alter this homology class. Transport along either
complex curve, away from its endpoints, preserves this positive
normal orientation. The two lifts belonging to each curve
therefore have the same sign.

Push (P3K.29.4) into the common exterior. It proves
\[
        m'_0=-m_0+2m_1+2m_2,\qquad
        m'_1=m_1,\qquad m'_2=m_2.                        \tag{P3K.29.5}
\]
The identification of common exteriors is the literal blowdown
isomorphism off the three curves; take the same compact
complement in that open set. The other two meridians can be
chosen away from both endpoint charts, where this map is the
identity. The new ambient orientation is the complex orientation,
so this open-set map is orientation preserving.

Thus the columns of the new absolute basis and of its relative
dual are, respectively,
\[
 M=
 \begin{pmatrix}-1&0&0\\2&1&0\\2&0&1\end{pmatrix},
 \qquad
 Q=M^{-T}=
 \begin{pmatrix}-1&2&2\\0&1&0\\0&0&1\end{pmatrix}.
                                                        \tag{P3K.29.6}
\]
In particular
\[
              D'_0=-D_0,\qquad
              D'_1=D_1+2D_0,\qquad
              D'_2=D_2+2D_0.                           \tag{P3K.29.7}
\]
All coefficients are integral, and \(M^2=I\). Relabelling
the selected edge gives the two other matrices recorded in
the exact certificate.

\par\medskip\noindent\textbf{3. The actual relative return and its variation.}\enspace

Use a fan retaining the physical delta edge and the actual
extending cusp product of K27. The positive parameter return
has a delta-equivariant representative \(F\) which is the
identity on a neighbourhood of the entire three-curve union.
To obtain it, take the closed Cox product on the cusp and
equivariant product connections on the remainder; their
partition-of-unity interpolation is still equivariant and is
fixed on the chosen cusp collar. This is the same physical
return, with its actual parameter closing.

The ordinary toric isomorphism, the common blowup charts and
the bounded Cox maps all intertwine the parameter products.
Hence a flop leaves this actual map \(F\) on the common
exterior unchanged. In particular no new relative boundary
framing is substituted in (P3K.29.6).

Since \(F\) is the identity near each \(C_i\), the marked
dual-cell generators show that it is the identity on
\(H_3(E_C)\) and \(H_4(E_C)\). Define its integral variation by
\[
 \operatorname{Var}_F:H_3(E_C,\partial E_C)\longrightarrow
 H_3(E_C),\qquad [D]\longmapsto[F(D)-D].
                                                        \tag{P3K.29.8}
\]
Let \(V\) be its matrix, with absolute basis \(m_i\) and
relative basis \(D_i\).

This matrix is symmetric. For a direct proof, double \(E_C\)
and extend \(F\) by the identity on the second half. The
absolute meridians and the doubles of \(D_i\) form the two
primitive middle-dimensional groups of the double. Their
intersection matrix has \(m_i\cdot D_j=\delta_{ij}\) and
\(D_j\cdot m_i=-\delta_{ij}\). The extended map fixes all
\(m_i\) and takes \(D_j\) to \(D_j+\sum_iV_{ij}m_i\).
Preservation of the intersections of two \(D\)'s gives
\(-V_{ij}+V_{ji}=0\). This proof does not require a choice
of a Lagrangian among the doubled \(D\)'s.

\par\medskip\noindent\textbf{4. The circle-trace kernel is an integral primitive lattice.}\enspace

The delta action is free on \(E_C\), since all its fixed points
are in the removed fixed curve \(C_0\). Put
\(B_C=E_C/S^1_\delta\). In the actual orbit \(Y=S^5\), the
image of \(C\) is the fixed two-sphere together with two
embedded arcs joining its two marked points. The two
nonfixed curves have free circle orbits away from their
endpoints, so their quotients are precisely these arcs.
The local Hopf chart supplies the regular-neighbourhood
map at the endpoints. Thus \(B_C\) is the complement of a
regular neighbourhood of this marked two-sphere-and-arcs
space in \(S^5\).

The same dual-cell construction gives
\[
 H_2(B_C)=\mathbb Z,\qquad H_3(B_C)=\mathbb Z^2,\qquad
 H_4(B_C)=0.
                                                        \tag{P3K.29.9}
\]
Both \(B_C\) and \(E_C\) are simply connected. The Euler
class \(x\) of the actual circle bundle is the primitive
generator of \(H^2(B_C)\). Its restriction to a meridian
two-sphere of the fixed curve is the original unit Hopf
bundle, with its given physical orientation.

Write \(\Delta\) for the circle-trace operation on relative
chains:
\[
 \Delta:H_2(E_C,\partial E_C)\longrightarrow
          H_3(E_C,\partial E_C).
                                                        \tag{P3K.29.10}
\]
It is induced by the actual action map
\(S^1\times E_C\to E_C\), with the positive delta circle.
The trace operation and the two Poincare--Lefschetz maps
identify (P3K.29.10), up to a single overall orientation
sign, with
\[
 H^4(E_C)\xrightarrow{\pi_!}H^3(B_C)
                    \xrightarrow{\pi^*}H^3(E_C).
                                                        \tag{P3K.29.11}
\]
One can see this equality directly on a transverse dual
chain: tracing its intersection once around the fibre is
the pullback of its oriented fibre integration. The
possibly opposite common sign does not change the image.

The integral Gysin sequence and (P3K.29.9) show that
\(\pi_!\) in (P3K.29.11) is an isomorphism. The pullback
\(\pi^*\) is injective. Its image is exactly the primitive
lattice of cohomology classes dual to \(m_1,m_2\).
Indeed \(m_0\) maps by the Hopf fibration to a two-sphere
and therefore has zero image in \(H_3(B_C)\). At an
interior point of either nonfixed curve, the quotient of
its circle tangent is a point, and its complex normal
four-disc maps diffeomorphically to a transverse normal
four-disc of the corresponding arc. Hence its meridian
maps, with degree of absolute value one, to the marked
arc meridian three-sphere. Those two meridians are the
primitive dual-cell basis of \(H_3(B_C)\). It follows
integrally that
\[
                       \operatorname{im}\Delta
                              =\mathbb ZD_1\oplus\mathbb ZD_2.
                                                        \tag{P3K.29.12}
\]

Equivariance of \(F\) gives the equality of actual chain
maps
\[
             \operatorname{Var}_3\Delta
                         =\Delta\operatorname{Var}_2.
                                                        \tag{P3K.29.13}
\]
The target of \(\operatorname{Var}_2\) is \(H_2(E_C)=0\).
Thus (P3K.29.12) is contained in the kernel of the middle
variation. Symmetry then gives
\[
                       V=\operatorname{diag}(v,0,0).
                                                        \tag{P3K.29.14}
\]
Filling the other two curve tubes and using the actual
relative-disc map in Section 1 identifies \(v\) with the
single fixed-curve variation in the orientations
(P3K.29.2). These signs differ from the ordered product
disc in K23, as follows explicitly. In its first-resolution
coordinates, the complex orientation is
\(D^3_y\times S^3_q\), and the positive complex meridian
is \(m_0=S^3_q\). Thus
\[
      m_0\cdot[D^3_y]=(-1)^{3\cdot3}=-1,\qquad
      D_0=-[D^3_y].
\]
K23's product variation \(t_{\rm K23}\) is the coefficient
of \(S^3_q\) in \(F(D^3_y)-D^3_y\), and it equals
\(-k_{\rm Jin}\). Consequently
\[
                v=-t_{\rm K23}=k_{\rm Jin}.
\]
Equivalently, the boundary of \(D_0\) has the complex
curve orientation, since \(y/|y|\) reverses that curve
orientation in K23. This is precisely the boundary-slope
orientation in (EVar.5), which gives the same integer.
For the primitive K22 families in K27,
\[
                              v=-3d.
                                                        \tag{P3K.29.15}
\]

\par\medskip\noindent\textbf{5. The matrix after the fixed-edge flop.}\enspace

The map \(F\) is the same common-exterior map, so its new
matrix is obtained by the two actual basis maps:
\[
 V'=M^{-1}VM^{-T}
   =v\begin{pmatrix}
       1&-2&-2\\
       -2&4&4\\
       -2&4&4
      \end{pmatrix}.                                   \tag{P3K.29.16}
\]
In particular the selected-edge variation remains \(v\),
and the two other diagonal entries become \(4v\).
These are variations of specified relative cycles.
They do not by themselves identify the relative stable
normal derivative of either of the latter edges.

For a flop of a nonfixed edge in a delta-retaining fan,
the relabelled matrix fixes the first meridian column.
Consequently the transformed variation is still
\(\operatorname{diag}(v,0,0)\). This follows by an exact
unimodular matrix calculation, not by discarding a cycle.

\par\medskip\noindent\textbf{6. Stable derivatives on the two original nonfixed cycles.}\enspace

There is also a direct integral vanishing for the original
two nonfixed cycles. Choose a stable frame of \(TB_C\)
by restricting a stable frame of the actual \(S^5\).
An invariant connection and the positive delta vector
give a stable frame of \(TE_C\). In this frame, the
derivative of the equivariant \(F\) is upper triangular:
it is the pullback of its quotient derivative on \(B_C\)
and the identity on the vertical line. The off-diagonal
connection term contracts linearly to zero, relative
to the boundary. Its stable derivative class is therefore
a pullback from \(B_C\).

The spaces are simply connected. The derivative and every
stable change of frame lift to the stable Spin group,
with the derivative equal to the identity on the boundary.
Its third homotopy generator is thus an integral
three-dimensional obstruction class, without a
two-to-one \(SO\)-Hurewicz ambiguity. A pullback class
evaluates to zero on the actual circle trace of a
two-cycle, because the composition of that trace with
\(\pi\) factors through the original two-cycle.
It therefore vanishes on \(D_1,D_2\).

This remains true in the single-curve product frame of
K23. Here is the exact frame-change formula. If
\(\sigma'=\sigma g\), then
\[
       c_{\sigma'}(F)(p)=
       g(Fp)^{-1}c_\sigma(F)(p)g(p).
                                                        \tag{P3K.29.17}
\]
The primitive third class of Spin, applied to a relative
three-cycle \(D\), changes by
\[
              -\langle g^*\eta,\operatorname{Var}_F(D)\rangle.
                                                        \tag{P3K.29.18}
\]
To verify this formula integrally, make \(g\) cellular and
constant on the two-skeleton, using the two-connectedness
of Spin. The product obstruction on each three-cell is
the sum of the three obstruction integers in
(P3K.29.17). The boundary terms cancel because \(F\) is
the identity there, leaving precisely the closed chain
\(F(D)-D\) with the minus sign in (P3K.29.18).

For \(D_1,D_2\), that variation is zero by
(P3K.29.14), so their stable derivative is zero in any
stable frame that extends to their single-curve exterior.
Choose the actual relative disc in Section 1. It misses
the other tubes, where the return is the identity, so the
same derivative evaluation, with the sign \(D_i=-[D^3_y]\)
just computed, is the negative of K23's full relative
single-curve derivative. Its value is therefore
\(+2n_{\rm rel}\), since K23 uses \(-2n_{\rm rel}\) on
the ordered product disc. The vanishing proves
\[
         n_{\rm rel}(C_1)=n_{\rm rel}(C_2)=0
             \quad\hbox{in a delta-retaining fan}.
                                                        \tag{P3K.29.19}
\]
This is an integer assertion in the specified actual
relative framing. It provides a second direct proof that
a nonfixed-edge flop in this fan sector leaves its
oriented homotopy-sphere class unchanged.

\par\medskip\noindent\textbf{7. Retained Pontryagin comparison.}\enspace

After the fixed-edge flop, the circle remains free on the
common exterior, but the new relative cycles are
(P3K.29.7). Their nonzero variations in
(P3K.29.16) make the frame term (P3K.29.18) essential.
The class \(p_1\) on the common exterior alone does not
give that frame term.

In the common blowup the two quotient maps contract the
two different rulings of \(E\). Their stable tangent
clutches differ on the corresponding normal meridian;
the local \(p_1\) Thom correction must be evaluated on
the actual doubled parameter cycle, including its
caps at the other two curves. Neither the static
meridian identity nor (P3K.29.16) assigns that correction.
Accordingly this statement does not assign the next
finite-isotropy edge's \(n_{\rm rel}\), its new central
class, or a further Jin pair.

The executable integer certificate is
\path{code/part3_fourfold/common_exterior_meridians.py} with its
JSON and captured output. It checks all three meridian
matrices, their dual relative matrices, the primitive
puncture relation, and every displayed variation change.
\end{proof}
